\section{Bode plots and stability margins}

\subsection{Bode magnitude and phase}

\subsubsection{Magnitude in decibels}
\[
|G(j\omega)|_{\mathrm{dB}}=20\log_{10}|G(j\omega)|
\]
Magnitude in decibels converts products of transfer-function factors into sums of asymptotic contributions.

\subsubsection{Phase angle}
\[
\phi(\omega)=\angle G(j\omega)
\]
The phase curve is the sum of phase contributions from gain, poles, zeros, integrators, and delays.

\subsubsection{Bode form}
\[
G(s)=K\frac{\prod_i(1+s/\omega_{z_i})}{s^q\prod_k(1+s/\omega_{p_k})}
\]
Break frequencies \(\omega_{z_i}\) and \(\omega_{p_k}\) determine where slopes and phases change.

\subsection{Asymptotic Bode rules}

\subsubsection{Real left-half-plane pole}
\[
\frac{1}{1+s/\omega_p}:\qquad
\Delta M=-20\ \mathrm{dB/dec},\quad
\Delta\phi=-90^\circ
\]
A stable real pole lowers the high-frequency magnitude slope and adds negative phase.

\subsubsection{Real left-half-plane zero}
\[
1+s/\omega_z:\qquad
\Delta M=+20\ \mathrm{dB/dec},\quad
\Delta\phi=+90^\circ
\]
A stable real zero raises the high-frequency magnitude slope and adds positive phase.

\subsubsection{Integrator}
\[
\frac{1}{s}:\qquad
M_{\mathrm{dB}}=-20\log_{10}\omega,\quad \phi=-90^\circ
\]
An integrator gives infinite DC gain and contributes \(-20\ \mathrm{dB/dec}\) at all frequencies.

\subsubsection{Right-half-plane zero}
\[
1-\frac{s}{\omega_z}:\qquad
\Delta M=+20\ \mathrm{dB/dec},\quad
\Delta\phi=-90^\circ
\]
A right-half-plane zero has the magnitude effect of a left-half-plane zero but the phase effect of a pole.

\subsubsection{Right-half-plane pole}
\[
\frac{1}{1-s/\omega_p}:\qquad
\Delta M=-20\ \mathrm{dB/dec},\quad
\Delta\phi=+90^\circ
\]
Right-half-plane poles make the open-loop plant unstable and require Nyquist reasoning for closed-loop stability.

\subsection{Second-order Bode factors}

\subsubsection{Second-order factor}
\[
G_2(s)=\frac{\omega_n^2}{s^2+2\zeta\omega_n s+\omega_n^2}
\]
At high frequency, a second-order pole pair contributes \(-40\ \mathrm{dB/dec}\) and \(-180^\circ\).

\subsubsection{Resonance peak}
\[
M_r=\frac{1}{2\zeta\sqrt{1-\zeta^2}}, \qquad 0<\zeta<\frac{1}{\sqrt{2}}
\]
A resonance peak indicates low damping and can create multiple crossover frequencies.

\subsection{Stability margins}

\subsubsection{Crossover frequency}
\[
|L(j\omega_c)|=1
\]
The crossover frequency is where the open-loop magnitude is \(0\ \mathrm{dB}\).

\subsubsection{Phase margin}
\[
\gamma_M=180^\circ+\angle L(j\omega_c)
\]
Positive phase margin indicates distance from the critical phase \(-180^\circ\) at crossover.

\subsubsection{Phase crossover frequency}
\[
\angle L(j\omega_\pi)=-180^\circ
\]
The phase crossover frequency is used to compute gain margin.

\subsubsection{Gain margin}
\[
A_M=\frac{1}{|L(j\omega_\pi)|},\qquad
A_{M,\mathrm{dB}}=-20\log_{10}|L(j\omega_\pi)|
\]
The gain margin is the multiplicative gain increase allowed before reaching the critical point.

\subsubsection{Bode stability criterion}
\[
\gamma_M>0,\qquad A_M>1
\]
For stable minimum-phase open-loop systems with a unique crossover, positive phase and gain margins indicate closed-loop stability.

\subsubsection{P-controller gain shift}
\[
L(s)=K_PG(s),\qquad
20\log_{10}|L(j\omega)|=20\log_{10}K_P+20\log_{10}|G(j\omega)|
\]
A positive proportional gain shifts only the magnitude curve and does not change the phase curve.
